\newif\ifhandoutmode
% \handoutmodetrue  %comment this line to enable animation

\ifhandoutmode
  \documentclass[handout, aspectratio=169,11pt]{beamer}
\else
  \documentclass[aspectratio=169,11pt]{beamer}
\fi

%──────────────────────────────────────────────────────────────────
%  Theme
%──────────────────────────────────────────────────────────────────
\usetheme{metropolis}
\metroset{progressbar=foot,subsectionpage=none}
\usepackage{appendixnumberbeamer}

%──────────────────────────────────────────────────────────────────
%  Palette  (semantic names — edit values here to retheme)
%──────────────────────────────────────────────────────────────────

% — Backgrounds (lightest → darkest) —
\definecolor{BgBase}{HTML}{F6F8FF}      % slide canvas
\definecolor{BgSurface}{HTML}{E6ECFF}   % block bodies, cards
\definecolor{BgElevated}{HTML}{CDD6FF}  % frametitle bar, raised tiles

% — Text —
\definecolor{TextPrimary}{HTML}{111827}   % body text
\definecolor{TextMuted}{HTML}{4B5563}     % captions, soft labels

% — Accent (indigo family) —
\definecolor{AccentMain}{HTML}{3D52CC}    % primary accent: titles, separators
\definecolor{AccentSoft}{HTML}{C7CFFF}    % tinted block headers, overlays

% — Warm highlight —
\definecolor{Highlight}{HTML}{BE185D}     % alerts, \important, \highlight
\definecolor{HighlightSoft}{HTML}{FBCFE8} % warm tint for decorative circles

%──────────────────────────────────────────────────────────────────
%  Beamer colour mapping
%──────────────────────────────────────────────────────────────────
\setbeamercolor{background canvas}  {bg=BgBase}
\setbeamercolor{normal text}        {fg=TextPrimary}
\setbeamercolor{frametitle}         {bg=BgElevated,   fg=TextPrimary}
\setbeamercolor{title separator}    {fg=AccentMain}
\setbeamercolor{progress bar}       {fg=AccentMain,   bg=BgElevated}
\setbeamercolor{progress bar in head/foot}{fg=AccentMain!80, bg=BgSurface}
\setbeamercolor{block title}        {bg=AccentSoft,   fg=TextPrimary}
\setbeamercolor{block body}         {bg=BgSurface,    fg=TextPrimary}
\setbeamercolor{alerted text}       {fg=Highlight}
\setbeamercolor{itemize item}       {fg=AccentMain}
\setbeamercolor{itemize subitem}    {fg=TextMuted}
\setbeamercolor{section in toc}     {fg=TextPrimary}

\makeatletter
\setlength{\metropolis@progressinheadfoot@linewidth}{2.2pt}
\makeatother

%──────────────────────────────────────────────────────────────────
%  Packages
%──────────────────────────────────────────────────────────────────
\usepackage{tikz}
\usetikzlibrary{arrows.meta,positioning,decorations.pathreplacing,
                shapes.geometric,backgrounds,calc,fit,fadings}
\usepackage{amsmath,amssymb,bm}
\usepackage{booktabs}
\newcommand{\faIcon}[1]{\textbullet}
\usepackage{xcolor}
\usepackage{graphicx}

\usepackage[
  backend=biber,
  style=numeric-comp,
  sorting=none,
  giveninits=true,
  maxbibnames=2,
  doi=false,
  url=false,
  eprint=false,
  isbn=false
]{biblatex}

\usetikzlibrary{overlay-beamer-styles}

\addbibresource{references.bib}

%──────────────────────────────────────────────────────────────────
%  Metadata
%──────────────────────────────────────────────────────────────────
\title{\textbf{Diffusion Models}}
\subtitle{How AI Learns to Create Images from Noise}
\author{Presented by: \textit{Anindya, Tamal and Din}}
\date{March 2, 2026}
\institute{CSE 200: Technical Writing and Presentation\\\textbf{Bangladesh University of Engineering and Technology}}

%──────────────────────────────────────────────────────────────────
%  Custom commands
%──────────────────────────────────────────────────────────────────
\newcommand{\emphasisword}[2]{\textcolor{#1}{\textbf{#2}}}
\newcommand{\highlight}[1]{\emphasisword{Highlight}{#1}}
\newcommand{\important}[1]{\emphasisword{Highlight}{#1}}
\newcommand{\soft}[1]{\textcolor{TextMuted}{#1}}

\hypersetup{
  bookmarksopen=true,
  bookmarksnumbered=true,
  pdfauthor={Anindya, Tamal and Din},
  pdftitle={Diffusion Models: How AI Learns to Create Images from Noise},
}

%══════════════════════════════════════════════════════════════════
\begin{document}
%══════════════════════════════════════════════════════════════════

%──────────────────────────────────────────────────────────────────
% TITLE SLIDE
%──────────────────────────────────────────────────────────────────
{
\setbeamercolor{background canvas}{bg=BgBase}

\begin{frame}
  \begin{tikzpicture}[remember picture,overlay]
    \foreach \r/\op in {5/0.30,3.5/0.40,2/0.60}{
      \fill[HighlightSoft,opacity=\op]
      (current page.north east) circle (\r cm);
      }
      \foreach \r/\op in {4/0.40,2.5/0.50,1.2/0.70}{
        \fill[AccentSoft,opacity=\op]
        (current page.south west) circle (\r cm);
        }
      \end{tikzpicture}
      \vspace{1.4cm}
      \titlepage
    \end{frame}
}

\begin{frame}{Outline}
  \begin{tikzpicture}[remember picture, overlay]
    \foreach \r/\op in {4.8/0.20,3.3/0.30,2.0/0.50}{
      \fill[HighlightSoft,opacity=\op]
      ([xshift=0.8cm]current page.east) circle (\r cm);
    }
    \foreach \r/\op in {3.8/0.30,2.4/0.40,1.3/0.60}{
      \fill[AccentSoft,opacity=\op]
      ([xshift=0.8cm,yshift=-1.2cm]current page.north east) circle (\r cm);
    }
  \end{tikzpicture}
  \tableofcontents
\end{frame}

%══════════════════════════════════════════════════════════════════
% SLIDE 1 — What Problem Are We Solving?
%══════════════════════════════════════════════════════════════════

\section{Introduction}

\begin{frame}{What Problem Are We Solving?}
  \vspace{4pt}
  \begin{figure}
    \centering
    \begin{minipage}[c]{\linewidth}
      \centering
      \begin{minipage}[b]{0.2\linewidth}
        \centering
        \includegraphics[width=\linewidth]{images/Real_Face.png}
        \\[1pt]
        {\tiny Real face}
      \end{minipage}
      \hspace{0.2\linewidth}
      \begin{minipage}[b]{0.2\linewidth}
        \centering
        \includegraphics[width=\linewidth]{images/AI_Generated_Face_05.png}
        \\[1pt]
        {\tiny AI-generated face \textit{(Sora)}}
      \end{minipage}
    \end{minipage}
  \end{figure}

  \vspace{4pt}
  \begin{itemize}
    \item Can a machine learn to \highlight{create new, realistic samples}
    that are \textit{indistinguishable} from real data?
    \pause
    \item In other words can it learn the underlying \important{probability distribution} of the data?
  \end{itemize}
  \pause
  \begin{center}
    \important{Diffusion models} provide a \textit{stable, high-quality} solution.
  \end{center}
\end{frame}

%══════════════════════════════════════════════════════════════════
% SLIDE 2 — The Core Idea: Learning by Destroying
%══════════════════════════════════════════════════════════════════
% \section{Core Idea}

\begin{frame}{The Core Idea: Learning by Destroying}

  \begin{center}
    \begin{tikzpicture}[
      stage/.style={rounded corners=6pt, draw=TextMuted!40,
                    minimum width=2.5cm, minimum height=2.4cm,
            font=\small, text=TextPrimary, align=center,
                    inner sep=6pt},
      bigarr/.style={-{Stealth[length=3mm]}, line width=1.6pt},
      forwardarr/.style={bigarr, draw=TextMuted},
      forwardlabel/.style={font=\footnotesize, text=TextMuted},
      reversearr/.style={bigarr, draw=Highlight},
      reverselabel/.style={font=\small, text=Highlight, align=center}]

      \path[use as bounding box] (-1.2,-2.8) rectangle (11.2,1.6);

      %% Forward (destroy)
      \node[visible on=<1->, stage, fill=AccentSoft] (clean) at (0,0)
        {\textbf{Clean Image}\\[6pt]\footnotesize contains\\all the detail};
      \node[visible on=<2->, stage, fill=BgElevated] (mid) at (5,0)
        {\textbf{Noisy Image}\\[6pt]\footnotesize some detail\\remains};
      \draw[visible on=<2->, forwardarr] (clean) -- (mid)
        node[midway, above, forwardlabel]{add noise};
      \node[visible on=<3->, stage, fill=HighlightSoft] (noise) at (10,0)
        {\textbf{Pure Noise}\\[6pt]\footnotesize no information\\left};
      \draw[visible on=<3->, forwardarr] (mid) -- (noise)
        node[midway, above, forwardlabel]{add noise};
      %% Reverse (learn)
      \draw[visible on=<4->, reversearr] (noise.south) -- ++(0,-0.5) -| (clean.south);
      \node[visible on=<4->, reverselabel] at (5,-2.1)
        {Learn to reverse this process step by step};

    \end{tikzpicture}
    \pause[5]

    \vspace{6pt}
    
    \begin{block}{We Essentially Have Two Diffusion Processes}
      \begin{itemize}
        \item The \textbf{Forward Diffusion Process} that gradually adds noise.\\
        \item The \textbf{Reverse Diffusion Process} that removes noise step by step.\\
      \end{itemize}
    \end{block}
  \end{center}
\end{frame}

%══════════════════════════════════════════════════════════════════
% SLIDE 3 — The Forward Process: Adding Noise
%══════════════════════════════════════════════════════════════════
% \section{Forward Diffusion Process}
\begin{frame}{The Forward Diffusion Process: Adding Noise}

  \vspace{10pt}
  \begin{center}
    \begin{minipage}{0.22\linewidth}
      \centering
      \includegraphics[height=3.0cm,keepaspectratio]{images/clean.png}\\[2pt]
      {\scriptsize Clean Image}
    \end{minipage}
    \pause
    \hfill
    \begin{minipage}{0.22\linewidth}
      \centering
      \includegraphics[height=3.0cm,keepaspectratio]{images/noised_slightly.png}\\[2pt]
      {\scriptsize Slightly Noisy}
    \end{minipage}
    \pause
    \hfill
    \begin{minipage}{0.22\linewidth}
      \centering
      \includegraphics[height=3.0cm,keepaspectratio]{images/noised_heavily.png}\\[2pt]
      {\scriptsize Heavily Noisy}
    \end{minipage}
    \pause
    \hfill
    \begin{minipage}{0.22\linewidth}
      \centering
      
\includegraphics[height=3.0cm,keepaspectratio]{images/noised_fully.png}\\[2pt]
      {\scriptsize Pure Noise}
    \end{minipage}
  \end{center}
  \pause

  \begin{block}{What Happens}
    Take a real image and add a \emph{tiny} bit of random noise.
    Repeat this hundreds of times until only noise remains.
  \end{block}
\end{frame}

%══════════════════════════════════════════════════════════════════
% SLIDE 4 — The Reverse Process: Removing Noise
%══════════════════════════════════════════════════════════════════
% \section{Reverse Diffusion Process}
\begin{frame}{The Reverse Diffusion Process: Removing Noise}

  \vspace{10pt}
  \begin{center}
    \begin{minipage}{0.22\linewidth}
      \centering
      
\includegraphics[height=3.0cm,keepaspectratio]{images/noised_fully.png}\\[2pt]
      {\scriptsize Pure Noise}
    \end{minipage}
    \pause
    \hfill
    \begin{minipage}{0.22\linewidth}
      \centering
      \includegraphics[height=3.0cm,keepaspectratio]{images/noised_heavily.png}\\[2pt]
      {\scriptsize Slightly Less Noisy}
    \end{minipage}
    \pause
    \hfill
    \begin{minipage}{0.22\linewidth}
      \centering
      \includegraphics[height=3.0cm,keepaspectratio]{images/noised_slightly.png}\\[2pt]
      {\scriptsize Almost Denoised}
    \end{minipage}
    \pause
    \hfill
    \begin{minipage}{0.22\linewidth}
      \centering
      \includegraphics[height=3.0cm,keepaspectratio]{images/clean.png}\\[2pt]
      {\scriptsize Fully Denoised}
    \end{minipage}
  \end{center}
  \pause

  \begin{block}{How Generation Works}
    \begin{itemize}
      \item Start from \important{pure noise}
      \item Predict and remove a little noise at each step
      \item Repeat until a realistic image appears
    \end{itemize}
  \end{block}
\end{frame}

\begin{frame}{Training}
  \begin{center}
    {\Large To produce realistic images, the model must \textbf{learn} to \important{reverse} the noise process.}\\[20pt]
    \pause
    {\setlength{\fboxsep}{8pt}\colorbox{BgElevated}{\Large \highlight{But how does the model learn this reverse process?}}}
  \end{center}
\end{frame}

%══════════════════════════════════════════════════════════════════
\section{Training}
% SLIDE 5 — How Do We Train It?
%══════════════════════════════════════════════════════════════════
\begin{frame}{How Do We Train It?}

  \begin{center}
    \begin{tikzpicture}[
      trainstep/.style={rounded corners=5pt, draw=TextMuted!30,
                   fill=BgSurface, text=TextPrimary,
                   minimum width=12.8cm, minimum height=1.0cm,
                   font=\large, inner sep=10pt, align=left},
      finalstep/.style={rounded corners=5pt, draw=AccentMain!40,
                   fill=BgElevated, text=TextPrimary,
                   minimum width=12.8cm, minimum height=1.0cm,
                   font=\large, inner sep=10pt, align=left},
                   arr/.style={->, TextMuted, line width=1.2pt}]
                   
      \node<1->[trainstep] (s1) at (0,0)
        {1.\; Take a real image from the training set};
      \node<2->[trainstep, below=14pt of s1] (s2)
        {2.\; Pick a random step and add the corresponding noise};
      \draw<2->[arr] (s1.south) -- (s2.north);
      \node<3->[trainstep, below=14pt of s2] (s3)
      {3.\; Ask the neural network: \textit{``what noise was added?''}};
      \draw<3->[arr] (s2.south) -- (s3.north);
      \node<4->[finalstep, below=14pt of s3] (s4)
        {4.\; Compare the prediction to the \textbf{actual noise}, minimize the error};
      \draw<4->[arr] (s3.south) -- (s4.north);
    \end{tikzpicture}
  \end{center}
\end{frame}

%══════════════════════════════════════════════════════════════════
% SLIDE 5.1 — NEW: The U-Net Architecture
%══════════════════════════════════════════════════════════════════
\begin{frame}{The ``Brain'' of the Model: The U-Net}
  \begin{center}
    \begin{tikzpicture}[
      box/.style={rounded corners=4pt, draw=AccentMain, fill=BgElevated, thick,
                  minimum width=2.8cm, minimum height=1.5cm, align=center},
      arrow/.style={-{Stealth[length=3mm]}, thick, TextMuted, line width=1.2pt}
    ]
      % Inputs
      \node[align=center] (input) at (0,0) {
        \framebox{\includegraphics[width=2cm,height=1.6cm]{images/noised_slightly.png}}\\[2pt]
        \scriptsize Noisy Image ($x_t$)
      };
      
      % Brain
      \node[box] (unet) at (5,0) {\textbf{U-Net}\\\textbf{Architecture}\\[4pt]\scriptsize Predicts Noise ($\epsilon_\theta$)};
      
      % Outputs
      \node[align=center] (output) at (10,0) {
        \framebox{
\includegraphics[width=2cm,height=1.6cm]{images/noised_fully.png}}\\[2pt]
        \scriptsize Predicted Noise ($\hat{\epsilon}$)
      };
      
      % Time step injection
      \node[align=center, text=Highlight] (time) at (5, 2) {\textbf{Time step ($t$)}};

      % Arrows
      \draw[arrow] (input) -- (unet);
      \draw[arrow] (unet) -- (output);
      \draw[arrow, Highlight, dashed] (time) -- (unet) node[midway, right] {\scriptsize injected};
    \end{tikzpicture}
  \end{center}
  
  \pause
  \vspace{10pt}
  \begin{itemize}
    \item \textbf{Input:} A noisy image and the specific time step ($t$).
    \item \textbf{Process:} The U-Net shrinks the image to understand the "global context," then expands it back, using \highlight{skip connections} to retain fine, pixel-level details.
    \item \textbf{Output:} An exact map of the noise predicted to have been added at step $t$.
  \end{itemize}
\end{frame}

%══════════════════════════════════════════════════════════════════
% SLIDE 6 — The Objective Function & Key Insight
%══════════════════════════════════════════════════════════════════
% \begin{frame}{The Objective Function \& Key Insight}
  
%   \begin{block}{The Simple Math Behind the Magic}
%     We train the model ($\epsilon_\theta$) to minimize the Mean Squared Error (MSE) between the \textbf{true noise} ($\epsilon$) and the \textbf{predicted noise}:
%     \vspace{-5pt}
%     \begin{equation*}
%       \mathcal{L} = \mathbb{E}_{t, x_0, \epsilon} \left[ \left\| \epsilon - \highlight{\epsilon_\theta(x_t, t)} \right\|^2 \right]
%     \end{equation*}
%   \end{block}
%   \pause

%   \vspace{10pt}
%   \begin{columns}[T]
%     \column{0.48\textwidth}
%       \begin{block}{Why This is Elegant}
%         \begin{itemize}
%           \item We control the \highlight{noise process}, so the target is mathematically defined.
%           \item Instead of predicting complex pixels directly, we predict simple noise. This turns generation into a precise \textbf{regression problem}.
%         \end{itemize}
%       \end{block}
      
%     \pause

%     \column{0.48\textwidth}
%       \begin{block}{Why It Works}
%         \begin{itemize}
%           \item The task is simplified into \textbf{small, predictable steps}.
%           \item Each step only removes a tiny amount of noise.
%           \item Repeating these simple steps reconstructs highly complex structures.
%         \end{itemize}
%       \end{block}
%   \end{columns}
% \end{frame}

% \section{Insights}
\begin{frame}{Key Insight \& Why It Works}
  \begin{columns}[T]
    \column{0.48\textwidth}
      \begin{block}{Key Insight}
        \begin{itemize}
          \item We control the \highlight{noise process}
          \item Added noise is mathematically defined
          \item So the training signal is precise
          \item This makes the objective a well-defined \textit{regression problem}
        \end{itemize}
      \end{block}
      
    \pause

    \column{0.48\textwidth}
      \begin{block}{Why It Works}
        \begin{itemize}
          \item The task is simplified into \textbf{small, predictable steps}
          \item Each step only removes a tiny amount of noise
          \item Repeating simple steps reconstructs complex structure
        \end{itemize}
      \end{block}
  \end{columns}
\end{frame}

%══════════════════════════════════════════════════════════════════
% SLIDE 6 — Why Diffusion Models Work So Well
%══════════════════════════════════════════════════════════════════

\begin{frame}{Why Diffusion Models Work So Well}
  \begin{block}{Advantages}
    \begin{itemize}
      \item \important{Stable training}
      \pause
      \item \important{High Quality Samples}
      \pause
      \item \highlight{Good Diversity} {\small (less mode collapse than GANs)}
      \pause
      \item \highlight{Flexible Conditioning} {\small (text-to-image, video-generation, audio synthesis)}
      \pause
      \item \highlight{Strong Theoretical Foundation} {\small (probabilistic modeling, score matching)}
    \end{itemize}
  \end{block}
  \pause

  \vspace{6pt}
  \begin{center}
    {In practice, diffusion models combine reliability with strong output quality.}
  \end{center}
\end{frame}

\begin{frame}{Summary}
  \begin{center}
    \Large
    Diffusion models generate realistic data by learning the \textit{underlying probability distribution} through reversing a \textit{gradual, structured noise process.}\\[20pt]
    \pause
    {\setlength{\fboxsep}{8pt}\colorbox{BgElevated}{\Large \highlight{But where is Diffusion Models actually applied?}}}
  \end{center}
\end{frame}

%══════════════════════════════════════════════════════════════════
% SLIDE 7 — Applications
\section{Applications}
%══════════════════════════════════════════════════════════════════
\begin{frame}{Applications}
  \vspace{6pt}
  \begin{center}
    \begin{tikzpicture}[
      app/.style={rounded corners=4pt,
                  minimum width=2.4cm, minimum height=0.58cm,
          font=\scriptsize\bfseries, text=TextPrimary,
                  inner sep=4pt, align=center},
      apparrow/.style={-{Stealth[length=2mm]},
                       TextMuted!65, line width=0.6pt},
      every node/.style={font=\scriptsize}]

      \node[draw=AccentMain, circle, fill=BgBase,
            font=\small\bfseries, text=AccentMain,
            minimum size=1.2cm] (c) at (0,0) {DMs};

      \node[visible on=<1->, app, fill=HighlightSoft]   (img)    at ( 90:2.4cm) {Image Synthesis};
      \draw[visible on=<1->, apparrow] (c) -- (img);
      \node[visible on=<2->, app, fill=BgElevated]      (aud)    at ( 30:2.4cm) {Audio / Music};
      \draw[visible on=<2->, apparrow] (c) -- (aud);
      \node[visible on=<3->, app, fill=HighlightSoft!60](vid)    at (330:2.4cm) {Video Generation};
      \draw[visible on=<3->, apparrow] (c) -- (vid);
      \node[visible on=<4->, app, fill=BgSurface]       (threed) at (270:2.4cm) {3-D / NeRF};
      \draw[visible on=<4->, apparrow] (c) -- (threed);
      \node[visible on=<5->, app, fill=AccentSoft]      (mol)    at (210:2.4cm) {Molecular Design};
      \draw[visible on=<5->, apparrow] (c) -- (mol);
      \node[visible on=<6->, app, fill=AccentSoft!60]   (med)    at (150:2.4cm) {Medical Imaging};
      \draw[visible on=<6->, apparrow] (c) -- (med);
    \end{tikzpicture}
  \end{center}
\end{frame}

\begin{frame}{Some Real-World Tools Powered by Diffusion}
  \begin{center}
    \begin{minipage}{0.36\textwidth}
      \centering
      
\includegraphics[height=1.0cm]{images/Dall-E-Logo-PNG.png}\\[1pt]
      \vspace{5pt}
      {\small \highlight{Text to Image (DALL·E)}}
    \end{minipage}
    \pause
    \vspace{20pt}
    
    \begin{minipage}{0.36\textwidth}
      \centering
      
\includegraphics[height=2cm]{images/Midjourney_Emblem_(in-colour).svg.png}\\
      {\small \highlight{Artistic Image Synthesis (Midjourney)}}
    \end{minipage}
    \pause
    \hspace{0.08\textwidth}
    \begin{minipage}{0.36\textwidth}
      \centering
      
\includegraphics[height=2cm]{images/OpenAI_Sora_icon.png}\\
      {\small \highlight{Text to Video (Sora)}}
    \end{minipage}
  \end{center}
\end{frame}

%══════════════════════════════════════════════════════════════════
% SLIDE 8 — Limitations & Open Challenges
%══════════════════════════════════════════════════════════════════
% \section{Limitations \& Challenges}

% \begin{frame}{Beyond Art: Science and Industry}
%   \begin{columns}[T]
%     \column{0.48\textwidth}
%       \begin{block}{Healthcare & Biology}
%         \begin{itemize}
%           \item \important{Protein Design:} Generating novel molecular structures for drug discovery (e.g., RFdiffusion).
%           \item \important{Medical Imaging:} Denoising and reconstructing high-resolution MRI and CT scans from sparse data.
%         \end{itemize}
%       \end{block}

%     \pause

%     \column{0.48\textwidth}
%       \begin{block}{Digital Media & Spatial}
%         \begin{itemize}
%           \item \highlight{Audio Synthesis:} Generating high-fidelity music and realistic sound effects from text.
%           \item \highlight{3D Generation:} Text-to-3D pipelines (e.g., DreamFusion) to accelerate game design and VR.
%         \end{itemize}
%       \end{block}
%     \end{columns}
%     \pause
% %   \vspace{15pt}
%   \begin{center}
%     {Diffusion is evolving from generating 2D pixel arrays to modeling complex, real-world data.}
%   \end{center}
% \end{frame}

%──────────────────────────────────────────────────────────────────
\begin{frame}{Current Limitations}
        \begin{itemize}
          \item \textbf{Slow generation:} requires many denoising steps
                (hundreds to thousands)
                \pause
          \item \textbf{High compute cost:} training requires significant
                GPU resources
                \pause
          \item \textbf{Large model size:} not easy to run on consumer hardware
        \end{itemize}
\end{frame}
\begin{frame}{Open Challenges}
        \begin{itemize}
          \item Can we make generation \highlight{faster}?
                (Active research area)
                \pause
          \item How do we ensure generated content is \important{safe and ethical}?
          \pause
          \item Can we extend to \highlight{longer videos} and
                \highlight{interactive content}?
        \end{itemize}
\end{frame}
% THANK YOU SLIDE
%──────────────────────────────────────────────────────────────────
\setbeamercolor{background canvas}{bg=BgBase}

% \section{Conclusion}

%──────────────────────────────────────────────────────────────────
% REFERENCES SLIDE
%──────────────────────────────────────────────────────────────────
\begin{frame}{References}
  \footnotesize
  \setlength{\bibitemsep}{4pt}
  \nocite{ho2020ddpm,song2019generative,song2020denoising,nichol2021improved}
  \printbibliography[heading=none]
\end{frame}

\begin{frame}
  \begin{tikzpicture}[remember picture,overlay]
    \foreach \r/\op in {6/0.20,4/0.30,2.2/0.50}{
      \fill[AccentSoft,opacity=\op]
        (current page.center) circle (\r cm);
    }
  \end{tikzpicture}
  \begin{center}
    {\Huge\bfseries\color{TextPrimary} Thank You!}\\[10pt]
    {\large\color{TextMuted} Thanks for listening!}
  \end{center}
\end{frame}

\end{document}